\documentclass{article}
\usepackage[english]{babel}
\usepackage{amsmath, amssymb}
\usepackage{enumitem}
\usepackage{geometry}
\geometry{margin=1in}

\title{STAT 17 Week 7 \\Central Limit Theorem (CLT) \& Sampling Distributions}
\author{Xiao Zhang}
\date{February 2026}

\begin{document}

\maketitle

\section*{1. Sampling Distributions}

\textbf{Population parameters} are fixed (but unknown).  
\textbf{Sample statistics} are random variables.

\begin{center}
\begin{tabular}{c|c}
\textbf{Population Parameter} & \textbf{Sample Statistic} \\ \hline
Mean $\mu$ & $\bar{x}$ \\
Std. Dev. $\sigma$ & $s$ \\
Proportion $p$ & $\hat{p}$ \\
\end{tabular}
\end{center}

\textbf{Definition.}  
A \textbf{sampling distribution} is the probability distribution of a statistic
(e.g., $\bar{x}$ or $\hat{p}$) from all possible samples of size $n$.

\bigskip

\section*{2. CLT for the Sample Mean}

Suppose:
\begin{itemize}[leftmargin=1.2cm]
\item Simple random sample (SRS)
\item Sample size $n$
\item Population mean $\mu$
\item Population standard deviation $\sigma$
\end{itemize}

\subsection*{Center}
\[
\mu_{\bar{x}} = \mu
\]
The sample mean is an unbiased estimator of $\mu$.

\subsection*{Spread (Standard Error)}
\[
\sigma_{\bar{x}} = \frac{\sigma}{\sqrt{n}}
\]
This is called the \textbf{standard error of the mean}.

As $n$ increases:
\[
\text{Standard Error decreases}
\]
Sample means vary less than individual observations.

\subsection*{Shape}

\textbf{Case 1: Population is normal}  
\[
\bar{x} \text{ is normal for any } n.
\]

\textbf{Case 2: Population not normal}

\textbf{Central Limit Theorem (CLT):}

If $n \ge 30$ (general guideline),
\[
\bar{x} \approx \text{Normal}
\]
regardless of the population shape.

Larger $n$ $\Rightarrow$ better approximation.

\subsection*{Full Model (When Conditions Met)}
\[
\bar{x} \sim N\left(\mu, \frac{\sigma}{\sqrt{n}}\right)
\]

\bigskip

\section*{3. Z-Score for Sample Means}

If sampling distribution is approximately normal:
\[
z = \frac{\bar{x} - \mu}{\sigma/\sqrt{n}}
\]

\textbf{Important distinction:}

\begin{itemize}
\item Individual observation:
\[
z = \frac{x - \mu}{\sigma}
\]
\item Sample mean:
\[
z = \frac{\bar{x} - \mu}{\sigma/\sqrt{n}}
\]
\end{itemize}

\bigskip

\section*{4. Conditions for CLT (Mean)}

Always check:

\begin{enumerate}
\item Random sample
\item Independence:
\[
n \le 0.05N
\]
\item Normality condition:
\begin{itemize}
\item Population normal OR
\item $n \ge 30$
\end{itemize}
\end{enumerate}

\bigskip

\section*{5. Sampling Distribution of the Sample Proportion}

Suppose:
\[
X \sim \text{Binomial}(n, p), \quad \hat{p} = \frac{X}{n}
\]

\subsection*{Center}
\[
\mu_{\hat{p}} = p
\]

\subsection*{Spread (Standard Error)}
\[
\sigma_{\hat{p}} =
\sqrt{\frac{p(1-p)}{n}}
\]

As $n$ increases:
\[
\text{Standard Error decreases}
\]

\subsection*{Shape Condition}

Sampling distribution of $\hat{p}$ is approximately normal if:
\[
np(1-p) \ge 10
\]

\subsection*{Full Model (When Conditions Met)}

If:
\begin{itemize}
\item Random sample
\item Independence
\item $np(1-p) \ge 10$
\end{itemize}

Then:
\[
\hat{p} \sim
N\left(p, \sqrt{\frac{p(1-p)}{n}}\right)
\]

\bigskip

\section*{6. Z-Score for Sample Proportions}

If approximately normal:
\[
z =
\frac{\hat{p} - p}
{\sqrt{\frac{p(1-p)}{n}}}
\]

Note:
\begin{itemize}
\item Use population $p$ in denominator.
\item Check independence: $n \le 0.05N$
\end{itemize}

\bigskip

\section*{7. Mean vs Proportion Summary}

\begin{center}
\begin{tabular}{c|c|c}
 & \textbf{Sample Mean} & \textbf{Sample Proportion} \\ \hline
Center & $\mu$ & $p$ \\
Spread & $\frac{\sigma}{\sqrt{n}}$ & $\sqrt{\frac{p(1-p)}{n}}$ \\
Normal Condition & $n \ge 30$ & $np(1-p) \ge 10$ \\
Z Formula & $\frac{\bar{x}-\mu}{\sigma/\sqrt{n}}$ 
& $\frac{\hat{p}-p}{\sqrt{p(1-p)/n}}$ \\
\end{tabular}
\end{center}


\end{document}
